\documentclass[10pt,a4paper]{amsart}
\usepackage{amsmath,amssymb,amsthm,graphicx}
% Retained preamble for the Stacks Project chapter snapshot in this directory.
% The chapter's own preamble.tex is a copy of the shared Stacks Project preamble; the block
% below is sliced MECHANICALLY, line for line, from it. Nothing was added, removed or
% redefined; the wrapper class, the xr-hyper/hyperref wiring and the externaldocument list
% are omitted because the assembler supplies the document class.
%
% lines taken from preamble.tex: 3, 5, 9, 11, 12, 156, 157, 158, ... (51 lines total)
%
% package coverage restored from the chapter's own preamble (2026-10-13): xy, xy-2cell

\usepackage[all]{xy}
\xyoption{2cell}
\UseAllTwocells

\newenvironment{reference}{\comment}{\endcomment}
\newenvironment{slogan}{\comment}{\endcomment}
\newenvironment{history}{\comment}{\endcomment}
\theoremstyle{plain}
\newtheorem{theorem}[subsection]{Theorem}
\newtheorem{proposition}[subsection]{Proposition}
\newtheorem{lemma}[subsection]{Lemma}
\theoremstyle{definition}
\newtheorem{definition}[subsection]{Definition}
\newtheorem{example}[subsection]{Example}
\newtheorem{exercise}[subsection]{Exercise}
\newtheorem{situation}[subsection]{Situation}
\theoremstyle{remark}
\newtheorem{remark}[subsection]{Remark}
\newtheorem{remarks}[subsection]{Remarks}
\numberwithin{equation}{subsection}
\def\lim{\mathop{\mathrm{lim}}\nolimits}
\def\colim{\mathop{\mathrm{colim}}\nolimits}
\def\Spec{\mathop{\mathrm{Spec}}}
\def\Hom{\mathop{\mathrm{Hom}}\nolimits}
\def\Ext{\mathop{\mathrm{Ext}}\nolimits}
\def\SheafHom{\mathop{\mathcal{H}\!\mathit{om}}\nolimits}
\def\SheafExt{\mathop{\mathcal{E}\!\mathit{xt}}\nolimits}
\def\Sch{\mathit{Sch}}
\def\Mor{\mathop{\mathrm{Mor}}\nolimits}
\def\Ob{\mathop{\mathrm{Ob}}\nolimits}
\def\Sh{\mathop{\mathit{Sh}}\nolimits}
\def\NL{\mathop{N\!L}\nolimits}
\def\CH{\mathop{\mathrm{CH}}\nolimits}
\def\proetale{{pro\text{-}\acute{e}tale}}
\def\etale{{\acute{e}tale}}
\def\QCoh{\mathit{QCoh}}
\def\Ker{\mathop{\mathrm{Ker}}}
\def\Im{\mathop{\mathrm{Im}}}
\def\Coker{\mathop{\mathrm{Coker}}}
\def\Coim{\mathop{\mathrm{Coim}}}
\DeclareMathSymbol{\boxtimes}{\mathbin}{AMSa}{"02}
\def\QCohstack{\mathcal{QC}\!\mathit{oh}}
\def\Cohstack{\mathcal{C}\!\mathit{oh}}
\def\Spacesstack{\mathcal{S}\!\mathit{paces}}
\def\Quotfunctor{\mathrm{Quot}}
\def\Hilbfunctor{\mathrm{Hilb}}
\def\Curvesstack{\mathcal{C}\!\mathit{urves}}
\def\Polarizedstack{\mathcal{P}\!\mathit{olarized}}
\def\Complexesstack{\mathcal{C}\!\mathit{omplexes}}
\def\Pic{\mathop{\mathrm{Pic}}\nolimits}
\def\Picardstack{\mathcal{P}\!\mathit{ic}}
\def\Picardfunctor{\mathrm{Pic}}
\def\Deformationcategory{\mathcal{D}\!\mathit{ef}}
\usepackage{cleveref}
\crefname{originallemma}{original lemma}{original lemmas}
\Crefname{originallemma}{Original lemma}{Original lemmas}
\makeatletter
\newlabel{cohomology-proposition-vanishing-Noetherian}{{02UZ}{source}{}{sourcepointer-1}{}}
\newlabel{cohomology-proposition-vanishing-Noetherian@cref}{{[originallemma][0][]02UZ}{[1][0][]source}{}{}{}}
\makeatother
\crefname{originallemma}{original lemma}{original lemmas}
\Crefname{originallemma}{Original lemma}{Original lemmas}
\makeatletter
\newlabel{lemma-pre-genus-curves}{{0E6I}{source}{}{sourcepointer-2}{}}
\newlabel{lemma-pre-genus-curves@cref}{{[originallemma][0][]0E6I}{[1][0][]source}{}{}{}}
\makeatother
\crefname{originallemma}{original lemma}{original lemmas}
\Crefname{originallemma}{Original lemma}{Original lemmas}
\makeatletter
\newlabel{more-algebra-lemma-lift-perfect-from-residue-field}{{0BCD}{source}{}{sourcepointer-3}{}}
\newlabel{more-algebra-lemma-lift-perfect-from-residue-field@cref}{{[originallemma][0][]0BCD}{[1][0][]source}{}{}{}}
\makeatother
\crefname{originallemma}{original lemma}{original lemmas}
\Crefname{originallemma}{Original lemma}{Original lemmas}
\makeatletter
\newlabel{spaces-perfect-lemma-flat-proper-perfect-direct-image-general}{{0CTM}{source}{}{sourcepointer-4}{}}
\newlabel{spaces-perfect-lemma-flat-proper-perfect-direct-image-general@cref}{{[originallemma][0][]0CTM}{[1][0][]source}{}{}{}}
\makeatother
\crefname{originallemma}{original lemma}{original lemmas}
\Crefname{originallemma}{Original lemma}{Original lemmas}
\makeatletter
\newlabel{spaces-perfect-lemma-proper-flat-h0}{{0E6B}{source}{}{sourcepointer-5}{}}
\newlabel{spaces-perfect-lemma-proper-flat-h0@cref}{{[originallemma][0][]0E6B}{[1][0][]source}{}{}{}}
\makeatother
\title{Stacks, Moduli of Curves: lemma-genus}
\author{The Stacks Project Authors}
\date{stacks-project:moduli-curves.tex:local-snapshot-20261012; GNU FDL 1.2 or later}
\begin{document}
\maketitle
\noindent\textit{Editorial note.} Statement span is lines 1127-1139 of sources/stacks-moduli-curves/pinned-original-source.tex; proof span is lines 1141-1165, the complete author proof (23 body lines).\par\medskip
\noindent\textit{Editorial note.} Navigation stubs: 5 cited-but-not-included labels are supplied as non-load-bearing pointers with their pinned tags where the snapshot tag index has them.\par\medskip
\noindent\textit{Editorial note.} extra\_packages is empty on purpose: the retained preamble already loads xy with options and a second declaration clashes in the assembler wrapper.\par\medskip
\noindent\textit{Editorial note.} Text is retained verbatim, including the author's own typos and placeholders.\par\medskip
\noindent\textit{Editorial note.} Queue index 21 of queue\_moduli-curves.json (pb=974, stmt 1127-1139, proof 1141-1165).\par\medskip
\noindent\textit{Allowed standard prerequisites (editorial).}\begin{itemize}
\item Explicit prerequisite boundary. The cited inputs are prerequisites; the argument in the retained span is not replaced by any cited result.
\end{itemize}\medskip
\ifdefined\hypertarget\hypertarget{sourcepointer-1}{}\fi
\noindent\textit{Source reference.} Original lemma 02UZ: sources/stacks-moduli-curves/pinned-original-source.tex, source lines 1165--1165. Cited by the retained statement or proof but not part of the collected spans; supplied only as a navigation pointer. The argument does not depend on its content.\par\medskip
\ifdefined\hypertarget\hypertarget{sourcepointer-2}{}\fi
\noindent\textit{Source reference.} Original lemma 0E6I: sources/stacks-moduli-curves/pinned-original-source.tex, source lines 1078--1088. Cited by the retained statement or proof but not part of the collected spans; supplied only as a navigation pointer. The argument does not depend on its content.\par\medskip
\ifdefined\hypertarget\hypertarget{sourcepointer-3}{}\fi
\noindent\textit{Source reference.} Original lemma 0BCD: sources/stacks-moduli-curves/pinned-original-source.tex, source lines 1165--1165. Cited by the retained statement or proof but not part of the collected spans; supplied only as a navigation pointer. The argument does not depend on its content.\par\medskip
\ifdefined\hypertarget\hypertarget{sourcepointer-4}{}\fi
\noindent\textit{Source reference.} Original lemma 0CTM: sources/stacks-moduli-curves/pinned-original-source.tex, source lines 1165--1165. Cited by the retained statement or proof but not part of the collected spans; supplied only as a navigation pointer. The argument does not depend on its content.\par\medskip
\ifdefined\hypertarget\hypertarget{sourcepointer-5}{}\fi
\noindent\textit{Source reference.} Original lemma 0E6B: sources/stacks-moduli-curves/pinned-original-source.tex, source lines 1165--1165. Cited by the retained statement or proof but not part of the collected spans; supplied only as a navigation pointer. The argument does not depend on its content.\par\medskip
\begin{lemma}
\label{lemma-genus}
Let $f : X \to S$ be a family of curves such that
$\kappa(s) = H^0(X_s, \mathcal{O}_{X_s})$ for all $s \in S$, i.e.,
the classifying morphism $S \to \Curvesstack$ factors
through $\Curvesstack^{h0, 1}$ (Lemma \ref{lemma-pre-genus-curves}). Then
\begin{enumerate}
\item $f_*\mathcal{O}_X = \mathcal{O}_S$ and this holds universally,
\item $R^1f_*\mathcal{O}_X$ is a finite locally free $\mathcal{O}_S$-module,
\item for any morphism $h : S' \to S$ if $f' : X' \to S'$ is the base change,
then $h^*(R^1f_*\mathcal{O}_X) = R^1f'_*\mathcal{O}_{X'}$.
\end{enumerate}
\end{lemma}

\begin{proof}
We apply Derived Categories of Spaces, Lemma
\ref{spaces-perfect-lemma-proper-flat-h0}.
This proves part (1). It also implies that locally on $S$
we can write $Rf_*\mathcal{O}_X = \mathcal{O}_S \oplus P$
where $P$ is perfect of tor amplitude in $[1, \infty)$.
Recall that formation of $Rf_*\mathcal{O}_X$ commutes
with arbitrary base change
(Derived Categories of Spaces, Lemma
\ref{spaces-perfect-lemma-flat-proper-perfect-direct-image-general}).
Thus for $s \in S$ we have
$$
H^i(P \otimes_{\mathcal{O}_S}^\mathbf{L} \kappa(s)) =
H^i(X_s, \mathcal{O}_{X_s})
\text{ for }i \geq 1
$$
This is zero unless $i = 1$ since $X_s$ is a $1$-dimensional
Noetherian scheme, see
Cohomology, Proposition \ref{cohomology-proposition-vanishing-Noetherian}.
Then $P = H^1(P)[-1]$ and $H^1(P)$ is finite locally free
for example by More on Algebra, Lemma
\ref{more-algebra-lemma-lift-perfect-from-residue-field}.
Since everything is compatible with base change we
also see that (3) holds.
\end{proof}

\end{document}
